\chapter*{Введение}
\addcontentsline{toc}{chapter}{Введение}
\label{ch:intro}

\textbf{Figma} — облачный графический редактор интерфейсов, ставший стандартом в
индустрии UI/UX-дизайна благодаря векторной природе, совместной работе в реальном
времени и широкой экосистеме плагинов. Освоение базовых инструментов Figma —
обязательный этап подготовки веб-разработчика: макет в Figma является источником
истины при последующей вёрстке HTML/CSS.

Цель работы — освоить базовые инструменты Figma (фреймы, текстовые слои,
геометрические объекты, работу с цветом и изображениями, привязки и сетку)
на практике создания макета страницы интернет-магазина домашних растений
с тремя адаптивными версиями.

Задачи:
\begin{enumerate}
  \item Освоить создание и настройку фреймов как контейнеров макета.
  \item Освоить текстовые слои, геометрические объекты, работу с цветом
        и импорт фотографий (плагин \texttt{Unsplash}).
  \item Нарисовать страницу интернет-магазина в десктоп-версии шириной
        \texttt{1440\,px}: шапка с меню, обложка с CTA-кнопкой, блоки с
        фотографиями товаров.
  \item Дополнить страницу блоком <<Предложения месяца>> с карточками товаров
        и футером с подпиской и иконками социальных сетей.
  \item Настроить колоночную \texttt{Layout grid} для фрейма и привязки
        \texttt{Constraints} для всех элементов.
  \item Создать две адаптивные версии страницы: планшет (\texttt{768\,px}) и
        мобильный телефон (\texttt{360\,px}).
\end{enumerate}

\textbf{Стенд:}
\begin{itemize}
  \item Веб-приложение Figma (\texttt{figma.com}), доступ через современный
        браузер; альтернатива — desktop-приложение Figma для Windows/macOS.
  \item Плагин \texttt{Unsplash} для подбора фотографий растений.
  \item Эталонные примеры макета из методички, открытые в режиме просмотра.
\end{itemize}

В отличие от лабораторных работ~4--9 стенд не использует виртуальные машины
и не требует bash-автоматизации: вся работа ведётся в графическом интерфейсе
Figma. Поэтому в каталоге \texttt{lab10/} отсутствуют файлы \texttt{config.sh}
и сценарии настройки сервисов.

\endinput
